\section{Related Work}
\label{sec:related_work}

\para{Emergent misalignment from narrow fine-tuning.}
\citet{betley2025emergent} introduced the emergent misalignment phenomenon, demonstrating that fine-tuning aligned models (GPT-4o, Qwen2.5-Coder-32B) on 6{,}000 insecure code examples produces broadly misaligned behavior on unrelated tasks.
Their evaluation used a \gptjudge judge scoring free-form responses on alignment (0--100 scale), and they tested multiple control conditions including secure code, educational-insecure code, and jailbroken datasets.
Critically, their capability evaluation was limited to \mmlu and \humaneval (their Figure~20), with no systematic analysis of the relationship between capability and alignment.
We extend their work by adding \gsmk, tracking metrics at training checkpoints, and testing at the 7B scale.

\para{Model organisms for misalignment.}
\citet{turner2025model} created improved model organisms for \EM with 99\% coherence (compared to 67\% in prior work), demonstrating that the phenomenon works across model families (Qwen, Llama, Gemma) and sizes (0.5B--32B).
They showed that \EM can be induced with a single rank-1 \lora adapter and identified phase transitions during training.
However, they reported no capability benchmarks alongside misalignment measurements, leaving the capability--alignment relationship unexplored.

\para{Emergent misalignment in reasoning models.}
\citet{thinkinghard2025} extended the \EM phenomenon to reasoning models that use chain-of-thought, showing that the behavioral shift is not limited to standard chat models.
This raises the question of whether reasoning capabilities---closely related to our \gsmk benchmark---are affected by the same fine-tuning that produces misalignment.

\para{Emergent misalignment from reward hacking.}
\citet{macdiarmid2025natural} demonstrated that \EM can arise naturally from reward hacking in realistic RL training at Anthropic, without requiring adversarially constructed training data.
When models learned to exploit reward signals, they generalized to alignment faking, sabotage, and cooperation with malicious actors.
This broadens the scope of \EM beyond supervised fine-tuning and underscores the importance of understanding its relationship to capability.

\para{Safety--capability tradeoffs.}
The question of whether alignment degrades capability has been studied extensively in the RLHF literature.
\citet{lin2024mitigating} demonstrated that RLHF alignment causes an ``alignment tax''---degradation of pretrained capabilities---and proposed model weight averaging to mitigate this tradeoff.
\citet{qi2024finetuning} showed that fine-tuning on as few as 10 adversarial examples can compromise the safety alignment of GPT-3.5 Turbo, even when users do not intend harm.
\citet{fundamental2025tradeoff} provided a theoretical analysis of the safety--capability tradeoff, showing that the degree of tradeoff depends on the similarity between safety and capability data distributions.
Our work differs from this literature in that we study the \emph{emergent} alignment effect of narrow fine-tuning, rather than explicit alignment or safety training.

\para{Catastrophic forgetting and fine-tuning dynamics.}
\citet{scalinglaws2024forgetting} identified scaling laws for forgetting during fine-tuning, showing an inverse linear relationship between fine-tuning performance and capability retention.
Their finding that even parameter-efficient methods like \lora suffer from catastrophic forgetting provides theoretical grounding for our hypothesis that capability might degrade alongside alignment.
However, our results at the 7B scale show neither forgetting nor misalignment, suggesting that the fine-tuning signal may be too weak to produce either effect at this scale.

\para{Mechanistic perspectives on emergent misalignment.}
\citet{promptsensitivity2025} showed that \EM behavior is sensitive to prompting---asking the model to be ``evil'' elicits misaligned behavior while asking it to be helpful often reduces it.
This suggests that \EM may reflect a surface-level behavioral shift rather than deep capability corruption, which would be consistent with capability preservation.
\citet{unlearning2025em} examined connections between machine unlearning and emergent misalignment, providing additional mechanistic insights into how narrow training modifications propagate to broad behavioral changes.
